\chapter{线性模型与近邻：先建立可靠基线}
\label{ch:linear}

线性模型把输入特征组合为一个得分；近邻模型根据相似训练记录预测。
二者都适合建立可解释的起点，但依赖不同假设。本章使用统一目标约定比较拟合、正则化和分类扩展。

\begin{keyidea}[把模型与训练目标配对]
连续数值可从平方损失线性回归开始；二分类概率可从逻辑回归开始。
特征表示决定“线性”能够表达什么，正则化强度与近邻数则需要训练内验证。
\end{keyidea}

\section{线性回归与最小二乘}
\label{sec:linear-ols}
线性回归（Linear Regression）预测$\hat y=b+\vct w^\top\vct x$。
普通最小二乘（Ordinary Least Squares, OLS）最小化残差平方和，
也可以等价地最小化其正比例缩放$J=\frac1{2n}\sum_i(\hat y_i-y_i)^2$。

\begin{figure}[htbp]
\centering\begin{tikzpicture}
\begin{axis}[width=124mm,height=55mm,xmin=-.2,xmax=2.3,ymin=.7,ymax=2.5,
 xlabel={$x$},ylabel={$y$},xtick={0,1,2},ytick={1,2},axis lines=left]
\addplot[Teal,thick,domain=-.1:2.2]{7/6+.5*x};
\addplot[only marks,mark=*,Amber,mark size=3pt] coordinates {(0,1)(1,2)(2,2)};
\draw[Quiet,dashed] (axis cs:0,1)--(axis cs:0,1.1666667);
\draw[Quiet,dashed] (axis cs:1,2)--(axis cs:1,1.6666667);
\draw[Quiet,dashed] (axis cs:2,2)--(axis cs:2,2.1666667);
\node[lab] at (axis cs:.65,2.3) {$\hat y=7/6+x/2$};
\end{axis}
\end{tikzpicture}

\caption{三个观测$(0,1),(1,2),(2,2)$的最小二乘直线；虚线为竖直残差。}
\par\smallskip
\readfigure{拟合得到$w=1/2$、$b=7/6$。三个预测为$7/6,5/3,13/6$，
按“真实减预测”计的残差为$-1/6,1/3,-1/6$，平方和为$1/6$，训练MSE为$1/18$。}
\end{figure}

\subsection{线性指的是系数的进入方式}
模型可以对$x^2$、对数或交互特征作线性组合，仍对系数线性。
线性回归本身不要求输入特征服从正态；若要作某些经典置信区间推断，才需额外的误差假设。

\quick{拟合关系不自动具有因果含义；外推到训练范围之外要结合领域结构判断。}

\section{怎样求参数：方程、分解与迭代}
\label{sec:linear-solve}
将截距作为一列1并记增广矩阵为$\widetilde X$，系数为$\vct\theta$。
令平方目标梯度为0，得到正规方程（Normal Equations）
\[
 \widetilde X^\top\widetilde X\vct\theta=\widetilde X^\top\vct y.
\]
只有当列满秩时，才能写$\hat\theta=(\widetilde X^\top\widetilde X)^{-1}\widetilde X^\top\vct y$。
不满秩时可能有多个最优系数；在样本上的拟合投影仍可确定。

\begin{center}
\begin{tabularx}{\textwidth}{P{35mm}Y}
\toprule
方法 & 实际考虑\\
\midrule
QR或SVD分解 & 通常比先形成正规方程再显式求逆更适合数值计算；SVD可处理秩亏。\\
梯度方法 & 避免显式矩阵求逆，适合较大或稀疏问题；需考虑尺度、步长和收敛。\\
闭式表达 & 有助于理解目标，但不是要求程序照着显式求逆执行。\\
\bottomrule
\end{tabularx}
\end{center}

正规方程不是“大数据时必定首选”或“绝不能用”的二元规则。
计算成本与样本数、特征数、稀疏性、条件数及求解器有关。
平方目标凸不代表任意步长的梯度下降都能收敛。

\begin{pitfall}[相关列与数值稳定性]
重复列或完整独热加截距可能造成秩亏；接近相关的列也可能使系数不稳定。
正则化和稳定求解都可以帮助，但用途与假设需要分别理解。
\end{pitfall}

\clearpage
\section{VIF：相关特征怎样放大系数方差？}
\label{sec:linear-vif}
多重共线性（Multicollinearity）指特征列之间存在较强线性关系。
即使预测相近，单个系数也可能对样本扰动敏感。方差膨胀因子
（Variance Inflation Factor, VIF）衡量经典OLS条件下这种\textbf{系数方差}的放大。

将第$j$列特征对其余特征和截距作辅助线性回归，记训练决定系数为$R_j^2$，则
\[
 \mathrm{VIF}_j=\frac1{1-R_j^2}.
\]
设含截距的设计矩阵列满秩，条件误差协方差为$\sigma^2 I$，令
$S_j=\sum_i(x_{ij}-\bar x_j)^2$，则
\[
 \operatorname{Var}(\hat w_j\mid X)=\frac{\sigma^2}{S_j}\,\mathrm{VIF}_j.
\]
比较基准保持$S_j$与误差方差相同；这里放大的是方差，标准误按$\sqrt{\mathrm{VIF}_j}$放大。

\begin{figure}[H]
\centering\begin{tikzpicture}
\begin{axis}[width=125mm,height=42mm,xmin=0,xmax=.95,ymin=1,ymax=20,
 ymode=log,log basis y=10,
 xlabel={辅助回归$R_j^2$},ylabel={VIF（对数轴）},xtick={0,.5,.8,.9},ytick={1,2,5,10,20},yticklabels={1,2,5,10,20}]
\addplot[Teal,thick,domain=0:.95,samples=121] {1/(1-x)};
\draw[Amber,dashed] (axis cs:.5,1)--(axis cs:.5,2)--(axis cs:0,2);
\addplot[only marks,mark=*,Amber] coordinates {(.5,2)};
\node[lab,anchor=west] at (axis cs:.14,5) {$R_j^2=0.5$，VIF=2};
\end{axis}
\end{tikzpicture}

\caption{辅助回归越接近完全解释该列，VIF越大；图上的0.5对应VIF=2。}
\end{figure}
\readfigure{取中心化特征$x_1=(-1,-1,1,1)^\top$、$x_2=(-1,0,0,1)^\top$。
内积为2、平方和为4与2，所以$R_1^2=2^2/(4\times2)=1/2$，VIF=2。
因此$\operatorname{Var}(\hat w_1\mid X)=\sigma^2/2$，而相同平方和的正交基准为$\sigma^2/4$。}

\begin{pitfall}[高VIF不等于“这个特征应立即删除”]
VIF不是系数数值的乘数，不直接衡量预测价值，更不识别因果关系。
保留、合并或正则化特征取决于解释还是预测的目标，应检查验证表现和领域含义。
5或10等经验界线不是普适科学定律；岭回归系数的方差也不能直接套用本页OLS公式。
\end{pitfall}

完全共线使$R_j^2=1$，此时普通系数可能无法唯一识别。
数值求解的条件数还受尺度影响；改变单位可改变系数方差，但在非退化情形下不改变相应VIF。


\section{多项式回归：改变表示，不增加样本}
\label{sec:linear-polynomial}
多项式回归（Polynomial Regression）使用$\phi(x)=(x,x^2,\ldots,x^d)$，
预测$\hat y=b+\sum_{j=1}^dw_jx^j$。对原输入可非线性，对系数仍线性。

\begin{center}
\begin{tabular}{lrrr}
\toprule
原输入$x$ & 截距列1 & $x$ & $x^2$\\
\midrule
0 & 1 & 0 & 0\\
1 & 1 & 1 & 1\\
2 & 1 & 2 & 4\\
\bottomrule
\end{tabular}
\end{center}
同样三条样本增加了表示列，没有新增独立观测。
前述三个观测的目标$[1,2,2]$可被$1+1.5x-0.5x^2$精确拟合，
但只有三个点不能证明这个二次关系适用于新数据。

\subsection{次数越高，越要检查数值与泛化}
高次幂可能尺度差异很大，并造成病态问题。可结合适当缩放、稳定基函数和正则化。
多变量展开还包含交互项，列数可能迅速增加。次数与表示步骤应在训练内选择，
不能看完测试结果再决定展开方式。

\quick{特征构造扩展了模型可表达的函数；它不是数据增强，也不是免费增加证据。}

\section{Ridge：连续收缩系数}
\label{sec:linear-ridge}
岭回归（Ridge Regression）在本章使用
\[
 J(\vct w,b)=\frac1{2n}\lVert X\vct w+b\vct1-\vct y\rVert_2^2
 +\frac\lambda2\lVert\vct w\rVert_2^2,\quad\lambda\geq0.
\]
截距$b$不惩罚。特征已中心化、单独处理截距时，
\[
 (X^\top X+n\lambda I)\hat{\vct w}=X^\top\vct y.
\]
$\lambda>0$时，对全部斜率施加的L2惩罚使此系数矩阵正定，即使原矩阵秩亏。

\begin{figure}[htbp]
\centering\begin{tikzpicture}
\begin{axis}[width=128mm,height=55mm,xmin=-2.5,xmax=2.5,ymin=-2.5,ymax=2.5,
 xlabel={无惩罚最优值 $z$},ylabel={惩罚后系数 $\hat w$},
 legend style={font=\small,at={(.02,.98)},anchor=north west},axis lines=left]
\addplot[Quiet,dashed,domain=-2.5:2.5]{x};\addlegendentry{无惩罚}
\addplot[Teal,thick,domain=-2.5:2.5]{x/1.5};\addlegendentry{Ridge，$\lambda=0.5$}
\addplot[Amber,thick] coordinates {(-2.5,-2)(-.5,0)(.5,0)(2.5,2)};\addlegendentry{Lasso，$\lambda=0.5$}
\end{axis}
\end{tikzpicture}

\caption{单系数教学目标$\frac12(w-z)^2$上，Ridge为$z/(1+\lambda)$；后文解释Lasso。}
\par\smallskip
\readfigure{Ridge通常使系数向0收缩，但不能说“绝不等于0”：当$z=0$时仍为0。
惩罚强度依赖目标中的$n$、$1/2$等缩放约定，比较软件参数前需对照完整公式。}
\end{figure}

尺度会改变同一系数惩罚所代表的限制，数值特征常在训练内标准化。
相关特征的单个系数路径可以变化，不能把每个系数都必然单调趋零当作一般定理。

\section{Lasso：为什么能够产生零系数？}
\label{sec:linear-lasso}
Lasso使用$J=\frac1{2n}\lVert X\vct w+b\vct1-\vct y\rVert^2+\lambda\lVert\vct w\rVert_1$，
截距不惩罚。L1惩罚在0处不可微，其几何与次梯度使零系数能够成为最优解。

\begin{figure}[htbp]
\centering\begin{tikzpicture}[x=13mm,y=13mm]
\begin{scope}
\draw[->,Quiet] (-1.5,0)--(1.8,0) node[right,lab] {$w_1$};
\draw[->,Quiet] (0,-1.4)--(0,1.7) node[above,lab] {$w_2$};
\filldraw[fill=Pale,draw=Teal,thick] (1,0)--(0,1)--(-1,0)--(0,-1)--cycle;
\node[lab] at (0,-2) {$|w_1|+|w_2|\leq1$};
\end{scope}
\begin{scope}[xshift=76mm]
\draw[->,Quiet] (-1.5,0)--(1.8,0) node[right,lab] {$w_1$};
\draw[->,Quiet] (0,-1.4)--(0,1.7) node[above,lab] {$w_2$};
\filldraw[fill=Pale,draw=Teal,thick] (0,0) circle (1);
\node[lab] at (0,-2) {$w_1^2+w_2^2\leq1$};
\end{scope}
\end{tikzpicture}

\caption{两个系数的L1与平方L2约束区域。这里只比较区域形状，未画任意相切等高线。}
\par\smallskip
\readfigure{L1菱形的角在坐标轴上，一个系数为0；L2圆盘边界更平滑。
这帮助理解稀疏性，但并不保证每个数据集都选择出真实变量。}
\end{figure}

对单系数目标$\frac12(w-z)^2+\lambda|w|$，精确解为软阈值
\[
 \hat w=\operatorname{sign}(z)\max(|z|-\lambda,0).
\]
当$z=0.3,\lambda=0.5$，Lasso系数为0，Ridge为0.2；
当$z=2$时，两者分别为1.5与$4/3$。

\begin{pitfall}[选中不等于因果；未选中不等于永远无用]
相关特征可能互相替代，选择结果也可能随样本变化。
Lasso在秩亏等情形下未必有唯一系数解。不能把L1也当成保证正规方程可逆的方法。
\end{pitfall}

\section{Elastic Net：同时控制稀疏与收缩}
\label{sec:linear-elastic}
弹性网络（Elastic Net）结合两种惩罚。本章约定总强度$\lambda$和混合比例$\rho$：
\[
 J=\frac1{2n}\lVert X\vct w+b\vct1-\vct y\rVert^2
 +\lambda\rho\lVert\vct w\rVert_1
 +\frac{\lambda(1-\rho)}2\lVert\vct w\rVert_2^2,
 \quad0\leq\rho\leq1.
\]
$\rho=1$得到Lasso，$\rho=0$得到本章约定的Ridge。
不同软件对alpha、lambda、l1-ratio的命名与缩放可能不同，应比较公式。

\begin{center}
\begin{tabularx}{\textwidth}{P{26mm}P{45mm}Y}
\toprule
方法 & 典型效果 & 选择时考虑\\
\midrule
Ridge & 连续收缩，通常保留多数非零系数。 & 相关列与稳定预测。\\
Lasso & 可以得到零系数。 & 稀疏表示，选择稳定性。\\
Elastic Net & 结合L1与L2，可同时稀疏与收缩。 & 相关特征组与两个超参数的权衡。\\
\bottomrule
\end{tabularx}
\end{center}

“特征很多就必须先用Elastic Net”不是通用规则。
应比较基线、特征结构、计算开销和验证表现。若$\lambda(1-\rho)>0$，
系数目标具有严格的二次惩罚，有助于唯一性；截距仍需按数据与模型定义处理。

\quick{正则化表达偏好，不会自动识别哪些特征是真实原因。}

\clearpage
\section{感知机：从分错的样本修正线性边界}
\label{sec:linear-perceptron}
感知机（Perceptron）使用得分$s=\vct w^\top\vct x+b$区分两个类别，
标签约定为$y\in\{-1,+1\}$。边界为$s=0$；这里预测时$s\geq0$判$+1$。
它直接学习分类边界，不输出已经校准的概率。

\begin{figure}[H]
\centering\begin{tikzpicture}
\begin{axis}[width=120mm,height=49mm,xmin=-2.7,xmax=2.7,ymin=-2.5,ymax=2.7,
 axis lines=middle,xlabel={$x_1$},ylabel={$x_2$},xtick={-2,-1,1,2},ytick={-2,-1,1,2},
 legend style={font=\scriptsize,at={(1.02,1)},anchor=north west,draw=none}]
\addplot[only marks,mark=*,Teal,mark size=2.7pt] coordinates {(2,1) (-1,2)};
\addlegendentry{$y=+1$}
\addplot[only marks,mark=triangle*,Amber,mark size=3pt] coordinates {(1,-1) (-2,-1)};
\addlegendentry{$y=-1$}
\addplot[Amber,dashed,thick,domain=-1.8:.7] {-2*x-1};
\addlegendentry{$2x_1+x_2+1=0$}
\addplot[Teal,thick,domain=-2.7:2.7] {-.5*x};
\addlegendentry{$x_1+2x_2=0$}
\end{axis}
\end{tikzpicture}

\caption{四个可分点。虚线是第一次更新后的边界，实线是第二次更新后的边界。}
\end{figure}
\readfigure{依次读取$(2,1),+1$与$(1,-1),-1$。
从$\vct w=(0,0),b=0$开始，步长1，第一次更新为$(2,1),1$；
第二条得分为2，符号错误，更新后为$(1,2),0$。此时四点都严格分对。}

本页采用在线规则：当$y_i(\vct w^\top\vct x_i+b)\leq0$时更新
\[
 \vct w\leftarrow\vct w+\eta y_i\vct x_i,\qquad
 b\leftarrow b+\eta y_i,\quad\eta>0;
\]
否则保持。边界点即使按预测平票规则判对，也按此训练规则更新。
可把输入增广为$(\vct x,1)$，将偏置一并理解为参数。

\begin{pitfall}[有限步结论需要可分条件]
对有界、以正间隔线性可分的固定训练数据，经典感知机规则有有限次更新的收敛保证。
非线性可分或含冲突标签时可能不停更新，需要训练上限或改用合适目标。
收敛只说明训练集分开，不保证最大间隔、概率校准或未来表现。
\end{pitfall}

阶跃输出几乎处处导数为0，不能把上述更新理解为直接穿过硬阶跃的普通梯度。
可从感知机损失$\max(0,-ys)$的次梯度解释更新；在$ys=0$处需选定约定。
逻辑回归使用平滑的逻辑损失，SVM使用间隔与惩罚；三者不能只因边界线性就混为同一算法。


\noindent\textbf{继续阅读：}Shalev-Shwartz与Ben-David，\textit{Understanding Machine Learning}（2014），第9章。


\section{逻辑回归：线性得分对应分类概率}
\label{sec:linear-logistic}
二元逻辑回归（Logistic Regression）令
$p=P(Y=1\mid\vct x)=\sigma(s)$，$s=b+\vct w^\top\vct x$，
等价于$\log[p/(1-p)]=s$，即对数几率对特征线性。

\begin{figure}[htbp]
\centering\begin{tikzpicture}[node distance=8mm]
\node[box,text width=27mm] (x) {输入$\vct x$};
\node[box,text width=34mm,right=of x] (s) {$s=b+\vct w^\top\vct x$\\实值得分};
\node[box,text width=30mm,right=of s] (p) {$p=\sigma(s)$\\预测概率};
\node[box,text width=26mm,right=of p] (y) {$p\geq t$？\\预测类别};
\draw[flow] (x)--(s);\draw[flow] (s)--(p);\draw[flow] (p)--(y);
\end{tikzpicture}

\caption{概率与类别之间还需要决策阈值$t$，默认0.5只是常见选择。}
\par\smallskip
\readfigure{若$s=\log3$，则几率为3，概率为$3/4$。
阈值0.5时预测1；阈值0.8时预测0。阈值改变决策，不改变已拟合概率函数。}
\end{figure}

以二元交叉熵拟合时，对单样本得分有$\partial\ell/\partial s=p-y$。
批量参数梯度为$X^\top(\vct p-\vct y)/n$，再加所选正则项。
对线性参数的逻辑损失为凸，但无正则且类别完全可分时，有限最优参数可能不存在。

\begin{pitfall}[线性得分不表示对输入的概率线性]
Sigmoid后的概率通常随输入非线性变化。
普通线性回归直接输出实数，可能超出$[0,1]$；逻辑回归为概率建模提供相应结构。
概率是否准确反映频率，还需校准与独立评估。
\end{pitfall}

\section{多分类：Softmax与一对其余}
\label{sec:linear-multiclass}
多项逻辑回归（Multinomial Logistic Regression）为$K$个互斥类别计算得分$s_k$，
令$p_k=\exp(s_k)/\sum_j\exp(s_j)$，用类别交叉熵联合拟合。
所有得分加同一常数不会改变概率，因此参数识别通常需基准类或其他约束。

\begin{center}
\begin{tabularx}{\textwidth}{P{28mm}YY}
\toprule
方法 & 怎样训练 & 输出如何理解\\
\midrule
Softmax & 一个联合的多分类目标。 & 类别概率和为1，面向互斥类别。\\
一对其余 OvR & 每类训练“此类/其余”的二元模型。 & 独立二元概率不必和为1；单标签预测需统一比较规则。\\
多标签 & 每个标签独立判断是否存在。 & 允许多个标签同时为真，不应强行Softmax归一化。\\
\bottomrule
\end{tabularx}
\end{center}

得分$[0,\log2,\log3]$经Softmax变为$[1/6,1/3,1/2]$。
三个OvR二元模型则可能同时输出$[0.7,0.6,0.2]$，它们并非一组已经归一化的互斥概率。
有些实现会另外归一化，但须按实现说明解释，不能与联合Softmax混为一谈。

\quick{先确定标签能否同时出现，再选择输出结构与损失。}

\section{\texorpdfstring{$k$}{k}近邻：根据局部训练记录预测}
\label{sec:linear-knn}
$k$近邻（$k$-Nearest Neighbors, KNN）查询最相似的$k$条训练记录。
分类可等权投票，回归可平均目标；也可按距离加权，需明确定义零距离的处理。

\begin{figure}[htbp]
\centering\begin{tikzpicture}[x=20mm,y=12mm]
\draw[->,Quiet] (-.3,0)--(5.7,0) node[right,lab] {$x$};
\foreach \x in {0,1}{\draw[Teal,thick,fill=Pale] (\x,0) circle (3pt);}
\foreach \x in {3,5}{\fill[Amber] (\x,0) circle (3pt);}
\foreach \x in {0,1,3,5}{\node[lab,below] at (\x,-.2) {$\x$};}
\node[draw=Ink,rectangle,fill=white,inner sep=2pt] (q) at (1.8,1.3) {?};
\draw[flow,dashed] (q)--(0,.1);\draw[flow,dashed] (q)--(1,.1);\draw[flow,dashed] (q)--(3,.1);
\node[lab,above] at (1.8,1.8) {查询$x=1.8$，$k=3$};
\node[lab] at (2.5,-1) {空心：类别0\qquad 实心：类别1};
\end{tikzpicture}

\caption{一维分类例子：训练位置0、1标签为0，3、5标签为1。}
\par\smallskip
\readfigure{查询1.8到四点距离为1.8、0.8、1.2、3.2。
三个近邻为1、3、0，票数2比1，预测0；若$k=1$也预测0，但依据只有最近点。}
\end{figure}

Minkowski距离为$d_p(\vct x,\vct z)=(\sum_j|x_j-z_j|^p)^{1/p}$，$p\geq1$。
$p=1$为曼哈顿，$p=2$为欧氏，$p\to\infty$为最大分量差。
尺度、相关列和无关特征都会影响“近”的含义，应根据任务设计表示。

\begin{pitfall}[奇数$k$并不消除所有平票]
二分类等权投票时奇数可避免两类票数相同；多分类、距离平局或加权规则仍可能平票。
KNN依赖局部相似性假设，也可能受异常点与高维影响。
大样本查询成本高，可考虑适合维度与精度需求的索引或近似检索。
\end{pitfall}

\begingroup\small
\noindent\textbf{参考：}ISLR第二版第3、4、6章；scikit-learn Linear Models与Nearest Neighbors文档。
本章例子和图独立构造，公式缩放均按正文约定。
\par\endgroup
